\section{系统模型与问题描述}
\label{rq3:sec:model}

\subsection{观测任务与传输反馈}
\label{rq3:subsec:tasks}
\label{rq3:subsec:channel}
\label{rq3:subsec:feedback}

考虑由感知节点集合 $\mathcal N=\{1,\ldots,N\}$ 与地面接收中心 $d$ 构成的协同观测系统，依次处理 $N_{\mathrm{task}}$ 个连续任务，其中第 $k$ 个任务的相对截止时间为 $D_k$。各节点的通信电量跨任务连续扣减，而接收状态与确认记录则在每个新任务开始时重置。

对于任务 $k$，上层内容选择模块依据各节点的局部观测与语义粒度，预先生成候选证据集合
\begin{equation}
  \mathcal{M}_k=\{m_{k,1},m_{k,2},\ldots,m_{k,K_k}\}.
  \label{rq3:eq:candidates}
\end{equation}
每条候选证据 $m$ 包含发送节点 $s(m)$、词元长度 $\ell_m$、表达粒度（标签、简洁声明、摘要或完整证据）以及对应的信念更新量。在交付控制阶段，候选内容、融合次序与融合规则均保持固定，控制器只需决定发送哪条候选，从而隐式确定对应的发送节点与传输开销。

记当前任务内决策时刻 $t$ 接收中心实际收到的唯一证据集合为 $S_{k,t}\subseteq\mathcal{M}_k$，则可能的接收状态空间为
\begin{equation}
  \states_k=2^{\mathcal{M}_k},\qquad |\states_k|=2^{K_k}.
  \label{rq3:eq:state-space}
\end{equation}
接收中心按既定次序对 $S_{k,t}$ 中的唯一证据执行融合，并输出预测类别 $c_k(S_{k,t})$。由于同一条证据仅融合一次，最终分类结果完全取决于接收集合本身，不受重复到达次数或到达先后顺序的影响。记该任务的参考类别为 $y_k$，集合 $S$ 对应的真实任务效用定义为
\begin{equation}
  u_k(S)=\ind\{S\ne\varnothing,\ c_k(S)=y_k\},
  \label{rq3:eq:true-utility}
\end{equation}
其中空集合效用为零，即未收到任何有效证据时不计作任务完成。需要说明的是，参考类别 $y_k$ 仅在离线训练与测试评价阶段可见，在线交付控制则采用第~\ref{rq3:subsec:utility} 节学习的效用网络进行估计。

在系统假设方面，考虑到单个观测任务的持续时间较短，假设任务内单跳无线接触保持连通、数据发送按串行方式执行、确认反馈占用独立控制信道，且各次发送的信道速率扰动相互独立。能耗统计涵盖数据帧与确认帧的无线发射支出，不计入候选元数据同步、端侧推理及无人机推进能耗。

底层链路沿用 LEADR 中的视距路径损耗、香农速率与发送能耗模型。由于数据串行传输且反馈信道独立，速率计算中不计同频并发数据干扰。结合节点到中心的距离 $d_{k,i}$、参考信道增益 $\rho_0$、有效带宽 $W_k$、噪声功率谱密度 $N_0$ 及数据发射功率 $P^{\mathrm d}$，候选 $m$ 所在链路的信道增益与预测速率分别为
\begin{equation}
  h_{k,i}=\frac{\rho_0}{\max\{d_{k,i},1\}^2},\qquad
  \bar r_{k,i}=W_k\log_2\left(1+\frac{P^{\mathrm d}h_{k,i}}{N_0W_k}\right).
  \label{rq3:eq:rate}
\end{equation}
上述预测速率在单个短任务内保持不变，距离与可用带宽则在任务间动态更新。

在物理层封装上，设平均每个 token 占 $b$ 字节，分组最大传输单元（Maximum Transmission Unit，MTU）的有效载荷上限为 $L_{\mathrm{MTU}}$ 字节，分组头部开销为 $H$ 字节，则候选 $m$ 实际需要传输的总比特数为
\begin{equation}
  B_m=8\left(b\ell_m+H\left\lceil\frac{b\ell_m}{L_{\mathrm{MTU}}}\right\rceil\right).
  \label{rq3:eq:bits}
\end{equation}
单次发送的预留时长记为
\begin{equation}
  \tau_{k,m}=\min\left\{\tau_{\max},\ \kappa\frac{B_m}{\bar r_{k,m}}\right\},
  \label{rq3:eq:duration}
\end{equation}
其中 $\kappa$ 为时间裕量系数，$\tau_{\max}$ 为单次发送时长上限，且整次传输必须落在任务剩余时间内。为刻画无线信道的随机波动，设第 $a$ 次尝试时的实际传输速率服从对数正态扰动：
\begin{equation}
  r_{k,m}^{(a)}=\bar r_{k,m}\exp\left(\sigma_r Z_a-\frac{\sigma_r^2}{2}\right),
  \qquad Z_a\sim\mathcal{N}(0,1),
  \label{rq3:eq:random-rate}
\end{equation}
其中 $\sigma_r$ 控制速率波动强度。若预留时长内能够传完全部 $B_m$ 比特，则该条证据完整送达；否则视为本次传输擦除，且依然耗尽完整预留时长。由此可得单次发送的成功概率：
\begin{equation}
  p_{k,m}=\Prb\{r_{k,m}^{(a)}\tau_{k,m}\ge B_m\}
  =\Phi\left(\frac{\ln(\bar r_{k,m}\tau_{k,m}/B_m)-\sigma_r^2/2}{\sigma_r}\right),
  \label{rq3:eq:success-probability}
\end{equation}
其中 $\Phi(\cdot)$ 为标准正态分布函数（$\sigma_r=0$ 时退化为确定性容量判定）。重传机制通过提供多次独立传输机会，提升证据至少成功送达一次的概率。

发送端依靠地面中心返回的累计确认（ACK）了解接收进展。如图~\ref{rq3:fig:delivery-scene} 所示，数据实际送达与发送端获知其送达往往并不同步：当确认尚在途中或已丢失时，继续重传可能造成不必要的重复占用，而盲目切换新候选又可能承担关键证据缺失的风险。

\begin{figure}[htbp]
\centering
% Illustrative events, without numerical timing or performance claims.
% Solid paths denote data; dashed paths denote cumulative feedback.
\begin{tikzpicture}[
  x=1mm,y=1mm,>=Stealth,
  font=\fontsize{9bp}{11.5bp}\selectfont,
  dline/.style={-{Stealth[length=1.6mm,width=1.2mm]},line width=.65pt,draw=blue!45!black},
  aline/.style={-{Stealth[length=1.6mm,width=1.2mm]},dashed,line width=.6pt,draw=orange!65!black},
  labelbg/.style={fill=white,inner sep=.4mm,align=center},
  pics/uav/.style={code={
    \draw[line width=.65pt] (-5.5,-2.5)--(5.5,2.5) (-5.5,2.5)--(5.5,-2.5);
    \foreach \x/\y in {-5.5/-2.5,-5.5/2.5,5.5/-2.5,5.5/2.5}{
      \draw[fill=black!3,line width=.45pt] (\x,\y) ellipse (3.4mm and 1.1mm);
      \draw[line width=.4pt] (\x-2.6,\y)--(\x+2.6,\y);
    }
    \draw[fill=black!12,line width=.5pt] (-2.2,-1.6) rectangle (2.2,1.6);
    \draw[fill=white,line width=.45pt] (-.9,-1.6) rectangle (.9,-3.5);
    \fill[black!65] (0,-2.8) circle (.45mm);
  }},
  pics/evidence/.style={code={
    \draw[fill=blue!4!white,line width=.4pt] (-4.5,-5.5) rectangle (4.5,5.5);
    \draw[line width=.35pt,draw=black!55] (-2.9,2.8)--(2.9,2.8)
      (-2.9,1.3)--(2.9,1.3) (-2.9,-.2)--(1.3,-.2);
  }}
]
\path[use as bounding box] (0,0) rectangle (148,-92);

% Two sensing UAVs and their candidate evidence.
\pic at (20,-16) {uav};
\pic at (68,-16) {uav};
\node at (20,-6) {感知无人机 A};
\node at (68,-6) {感知无人机 B};
\pic at (18,-32) {evidence};
\node at (18,-35) {$m_A$};
\pic at (65,-32) {evidence};
\node at (65,-35) {$m_B$};
\draw[draw=black!45,line width=.35pt] (19,-21)--(18,-26.5);
\draw[draw=black!45,line width=.35pt] (67,-21)--(65,-26.5);
\node at (18,-42) {已发送证据};
\node at (65,-42) {待发候选};

% Ground receiver: antenna and edge terminal, at natural scene proportions.
\draw[fill=black!2,line width=.45pt] (114,-49) rectangle (146,-52);
\draw[line width=.6pt] (124,-31)--(119,-49)--(129,-49)--cycle;
\draw[line width=.4pt] (122,-38)--(126,-38) (121,-43)--(127,-43);
\draw[line width=.4pt] (120,-47)--(126,-38) (128,-47)--(122,-38);
\fill[black!65] (124,-31) circle (.8mm);
\draw[line width=.45pt] (121.5,-28) arc[start angle=140,end angle=40,radius=3.2mm];
\draw[line width=.4pt] (119.8,-26.6) arc[start angle=140,end angle=40,radius=5.4mm];
\draw[fill=black!5,line width=.5pt] (132,-40) rectangle (143,-46.5);
\draw[line width=.5pt] (137.5,-46.5)--(137.5,-49);
\draw[line width=.5pt] (134,-49)--(141,-49);
\node at (129.5,-56) {地面接收中心};
\draw[draw=black!35,line width=.4pt] (3,-53)--(146,-53);

% An already delivered message and a prospective alternative.
\draw[dline] (29,-18.5) to[out=-7,in=166] (119,-34);
\node[labelbg] at (84,-28) {$m_A$ 已到达};
\draw[dline,draw=green!40!black] (77,-18) to[out=0,in=155] (118,-27);
\node[labelbg] at (101,-16) {可新发 $m_B$};
\draw[aline] (116,-46.5) -- (36,-46.5) to[out=180,in=-40] (28,-23);
\node[labelbg] at (80,-48) {累计确认延迟或丢失};

% A short event view separates sender knowledge from receiver reality.
\draw[draw=black!25,line width=.4pt] (0,-60)--(148,-60);
\node[anchor=east] at (24,-68) {发送端};
\node[anchor=east] at (24,-81) {接收端};
\draw[->,draw=black!55,line width=.4pt] (27,-68)--(146,-68);
\draw[->,draw=black!55,line width=.4pt] (27,-81)--(146,-81);
\node[anchor=north east] at (146,-85) {时间};
\draw[dline] (36,-68)--(57,-81);
\fill[blue!45!black] (36,-68) circle (.6mm);
\fill[blue!45!black] (57,-81) circle (.6mm);
\node[labelbg] at (45,-63.5) {发送 $m_A$};
\node[labelbg] at (52,-87) {$m_A$ 已接收};
\draw[aline] (57,-81)--(86,-69.5);
\draw[draw=orange!65!black,line width=.7pt] (84.8,-70.7)--(87.2,-68.3)
  (84.8,-68.3)--(87.2,-70.7);
\node[labelbg] at (74,-77) {确认未到};
\draw[draw=black!50,line width=.4pt] (96,-62)--(96,-88.5);
\node[labelbg] at (96,-63.5) {当前};
\node[labelbg] at (121,-73) {补发 $m_A$ 或新发 $m_B$};
\node[labelbg] at (120,-86) {实际已收 $\{m_A\}$\\重复不新增证据};
\end{tikzpicture}

\bicaption{语义数据送达与累计确认不同步的通信示意}{Semantic data delivery with delayed or missing cumulative acknowledgment}
\label{rq3:fig:delivery-scene}
\end{figure}

具体而言，每次数据发送结束后，接收中心均会生成一条包含尝试序号 $j$ 与当前接收集合快照 $S_j^{\mathrm{ack}}$ 的累计确认帧（短控制头可独立标识调度序号，因此即便本次数据载荷解码失败，中心仍可回传最新接收位图，以便发送端同步确认已收集合与漏收情况）。记第 $j$ 次数据发送的结束时刻为 $t_j^{\mathrm{end}}$，反馈帧长度为 $B^{\mathrm f}$ 比特，则各确认帧沿对应链路按预测速率串行发出的完成时刻满足
\begin{equation}
  v_j=\max\{t_j^{\mathrm{end}},v_{j-1}\}+\tau_j^{\mathrm f},
  \qquad \tau_j^{\mathrm f}=\frac{B^{\mathrm f}}{\bar r_{k,m_j}},
  \label{rq3:eq:feedback-serialization}
\end{equation}
其中初始值 $v_{-1}=0$。串行发出后，每条确认帧以概率 $\epsilon_{\mathrm f}$ 发生丢失；若未丢失，则经历随机额外时延 $L_j$ 到达发送端：
\begin{equation}
  L_j\sim\mathcal{U}\left(\frac{\bar d_{\mathrm f}}{2},\frac{3\bar d_{\mathrm f}}{2}\right),
  \qquad t_j^{\mathrm{ack}}=v_j+L_j.
  \label{rq3:eq:feedback-delay}
\end{equation}
数据成败、反馈丢失及反馈延迟彼此独立抽样。由于确认帧可能乱序到达，控制器仅保留尝试序号最大的有效快照，不会用陈旧位图覆盖较新状态。

记时刻 $t$ 发送端掌握的全部可见历史为 $\history_{k,t}$，涵盖候选与链路预测、过往发送记录、已到达的累计确认以及当前剩余时间与电量（不包含真实的隐藏接收集合和在途确认）。定义接收信念为当前真实接收集合的条件概率分布：
\begin{equation}
  b_{k,t}(S)=\Prb\{S_{k,t}=S\mid\history_{k,t}\},\qquad S\in\states_k.
  \label{rq3:eq:belief-definition}
\end{equation}
在数据成败与反馈丢失相互独立的前提下，即便尚未收到确认，发送端仍可结合历史发送记录推断各接收状态的概率，而主动等待能否换取新状态信息则取决于在途确认的到达概率。

\subsection{通信资源与优化目标}
\label{rq3:subsec:energy}
\label{rq3:subsec:objective}

设感知节点 $i$ 的初始通信电量为 $E_i^{(0)}$，则单次数据发送与对应确认回传的发射能耗分别为
\begin{equation}
  e_{k,m}^{\mathrm d}=P^{\mathrm d}\tau_{k,m},\qquad
  e_{k,m}^{\mathrm f}=P^{\mathrm f}\tau_{k,m}^{\mathrm f}.
  \label{rq3:eq:transmit-energy}
\end{equation}
其中数据发送能耗从对应感知节点 $s(m)$ 电量中逐次扣减；接收中心独立供电，但反馈帧一经生成（无论后续是否丢失或超时到达），其发射能耗 $e_{k,m}^{\mathrm f}$ 均计入系统总通信成本。记节点 $i$ 在第 $k$ 项任务中累计消耗的数据发送能量为 $e_i(k)$，则跨任务尺度的节点剩余电量演化为
\begin{equation}
  E_i(k+1)=E_i(k)-e_i(k),\qquad E_i(1)=E_i^{(0)}.
  \label{rq3:eq:energy-evolution}
\end{equation}

在每个决策时刻，控制器可执行三类动作：发送某条候选 $m$（若该候选此前已发过即为重传）、主动等待时长 $w$ 或直接停止当前任务。记当前已执行尝试次数为 $n_{k,t}$，单任务尝试上限为 $A_{\max}$，则同时满足剩余时限、节点电量与尝试上限的可行发送集合为
\begin{equation}
  \begin{split}
  \actions^{\mathrm{send}}_{k,t}=\{m\in\mathcal{M}_k:\;&
  \tau_{k,m}\le D_k-t,\\[-2pt]
  &e_{k,m}^{\mathrm d}\le E_{s(m)}(k,t),\quad n_{k,t}<A_{\max}\}.
  \end{split}
  \label{rq3:eq:action-mask}
\end{equation}
等待动作须满足 $0<w\le D_k-t$，停止动作则始终可行。通过式~\eqref{rq3:eq:action-mask} 的硬约束校验，系统严格杜绝超时发送或电量透支。

记策略 $\pi$ 在任务 $k$ 上的停止时刻为 $T_k$、终止时的真实接收集合为 $S_k^{\mathrm{end}}$、实际执行的数据发送尝试集合为 $\mathcal{T}_k$，则该任务的总通信能耗为
\begin{equation}
  E_k^{\mathrm{comm}}=\sum_{a\in\mathcal{T}_k}
    \left(e_{k,m_a}^{\mathrm d}+e_{k,m_a}^{\mathrm f}\right).
  \label{rq3:eq:total-communication-energy}
\end{equation}
在有限任务序列中，交付控制的目标是在提升按时正确完成率与压缩通信能耗之间取得平衡，相应的全局随机优化问题可表述为
\begin{equation}
  \begin{aligned}
    \max_{\pi}\quad &\frac{1}{N_{\mathrm{task}}}\E_{\pi}\left[
      \sum_{k=1}^{N_{\mathrm{task}}}\left(u_k(S_k^{\mathrm{end}})-\lambda_E
      \frac{E_k^{\mathrm{comm}}}{E_{\mathrm{ref}}}\right)\right] \\
    \text{s.t.}\quad &\sum_{k=1}^{N_{\mathrm{task}}} e_i(k)\le E_i^{(0)},\quad i\in\mathcal{N},\\
    &0\le T_k\le D_k,\quad |\mathcal{T}_k|\le A_{\max},\\
    &a_{k,t}\text{ 仅依赖 }\history_{k,t}\text{，且满足式~\eqref{rq3:eq:action-mask}},
  \end{aligned}
  \label{rq3:eq:global-problem}
\end{equation}
其中 $E_{\mathrm{ref}}>0$ 为能量归一化尺度，$\lambda_E\ge0$ 为能耗惩罚系数。该问题的难点在于接收状态隐藏、真实任务效用不可观测以及多任务共享有限电量三者交织：多发一次数据既可能补齐当前判断证据，也会直接压缩后续任务的能量余量。为此，下节先利用离线学习的效用网络替代不可观测的真实效用，再基于接收信念、跨任务虚拟能量队列与有限前瞻，构造式~\eqref{rq3:eq:global-problem} 的在线近似求解方法。本章主要符号汇总于表~\ref{rq3:tab:notation}。

\begin{table}[t]
  \centering\small
  \bicaption{语义交付控制模型的主要符号}{Main notation of the semantic delivery control model}
  \label{rq3:tab:notation}
  \begin{tabularx}{\linewidth}{@{}lY@{}}
    \toprule
    符号 & 含义\\
    \midrule
    $N,\ N_{\mathrm{task}}$ & 感知节点数及连续执行的任务总数\\
    $\mathcal{M}_k,\ K_k$ & 任务候选证据集合及其规模\\
    $S_{k,t},\ b_{k,t}(S)$ & 实际接收集合及其条件概率\\
    $c_k(S),\ u_k(S)$ & 接收集合对应的预测类别与真实任务效用\\
    $\hat q_\theta(S),\ q_k(S)$ & 原始预测效用与一致性修正效用\\
    $\bar r_{k,m},\ p_{k,m}$ & 预测链路速率与单次送达概率\\
    $\tau_{k,m},\ D_k$ & 数据发送时长与任务截止时间\\
    $\bar d_{\mathrm f},\epsilon_{\mathrm f}$ & 反馈平均额外时延与反馈丢失概率\\
    $E_i(k,t),\ Z_i(k)$ & 节点剩余发送能量与虚拟能量队列\\
    $\pi_{k,m}$ & 一次发送及其反馈的归一化资源价格\\
    $\eta_E,\eta_T$ & 局部控制的能量权重与时间权重\\
    \bottomrule
  \end{tabularx}
\end{table}
\FloatBarrier